% =====================================================================
%  Documento independiente para la Dra. Kruse (21 oct. 2026): las secciones
%  que todavia no reviso (1 Introduccion, 2 Estado del arte, 3 Marco teorico)
%  y las nuevas (7 Perspectivas futuras, 8 Conclusiones), con sus referencias.
%  Compilar con:  latexmk introduccion_conclusiones_sol.tex
%  (requiere haber compilado main.tex antes: toma de build/main.aux los numeros
%  de las referencias cruzadas a otras secciones).
% =====================================================================
\documentclass[12pt,a4paper,oneside]{report}
% =====================================================================
%  preambulo.tex — Configuración general del PFC
%  Compilar con: latexmk  (pdflatex + biber; ver .latexmkrc)
% =====================================================================

% ---------- Idioma y codificación ----------
\usepackage[T1]{fontenc}
\usepackage[utf8]{inputenc}
\usepackage[english,spanish,es-tabla,es-noshorthands,es-nodecimaldot]{babel} % español = idioma principal (último)
\usepackage{csquotes}

% ---------- Tipografía (la plantilla ITBA usa Noto Serif) ----------
\usepackage{microtype}
\usepackage{newunicodechar}
% Red de seguridad por si se pega texto con caracteres Unicode sueltos
\newunicodechar{µ}{\textmu}
\newunicodechar{−}{\ensuremath{-}}
\newunicodechar{≤}{\ensuremath{\leq}}
\newunicodechar{≥}{\ensuremath{\geq}}
\newunicodechar{⁺}{\ensuremath{^{+}}}
\newunicodechar{→}{\ensuremath{\rightarrow}}
\newunicodechar{η}{\ensuremath{\eta}}
\newunicodechar{λ}{\ensuremath{\lambda}}

% ---------- Página ----------
\usepackage[a4paper,top=2.5cm,bottom=2.5cm,left=2.5cm,right=2.5cm]{geometry}
\usepackage{setspace}
\onehalfspacing
\setlength{\parskip}{0.5em}
\setlength{\parindent}{0pt}
\setlength{\emergencystretch}{2em}  % evita renglones desbordados

% ---------- Matemática, unidades, tablas, figuras ----------
\usepackage{amsmath}
\usepackage{notomath}   % Noto Serif para texto y matemática (coherente con la plantilla)  % cargar después de amsmath; ya provee los símbolos de amssymb
\usepackage[locale=DE,per-mode=symbol]{siunitx} % DE → coma decimal
\usepackage{graphicx}
\usepackage{tikz}
\usepackage{booktabs}
\usepackage{tabularx}
\usepackage{array}
\usepackage{multirow}
\usepackage{pifont}          % \ding{173} = ②
\usepackage[font=small,labelfont=bf,labelsep=period]{caption}
\usepackage{subcaption}
\renewcommand{\thesubfigure}{\Alph{subfigure}} % paneles (A), (B), ... igual que en los pies de figura
\usepackage{placeins} % \FloatBarrier: evita que las figuras de una seccion floten a la siguiente
\usepackage{float}    % [H]: figura exactamente donde esta en el texto (se usa en el cap. 4)
\usepackage{enumitem}
\usepackage[dvipsnames]{xcolor}
\newcolumntype{C}{>{\centering\arraybackslash}X}
\graphicspath{{figuras/}{logos/}}

% ---------- Títulos: "1. Introducción", "1.1. ..." (sin "Capítulo") ----------
\usepackage{titlesec}
\titleformat{\chapter}[hang]{\normalfont\LARGE\bfseries}{\thechapter.}{0.5em}{}
\titlespacing*{\chapter}{0pt}{0pt}{1.2em}
\titleformat{\section}{\normalfont\Large\bfseries}{\thesection.}{0.5em}{}
\titleformat{\subsection}{\normalfont\large\bfseries}{\thesubsection.}{0.5em}{}
\setcounter{secnumdepth}{3}
\setcounter{tocdepth}{2}

% ---------- Encabezados y pie ----------
\usepackage{fancyhdr}
\pagestyle{fancy}
\fancyhf{}
\fancyfoot[C]{\thepage}
\renewcommand{\headrulewidth}{0pt}
\fancypagestyle{plain}{\fancyhf{}\fancyfoot[C]{\thepage}\renewcommand{\headrulewidth}{0pt}}

% ---------- Bibliografía IEEE ----------
% language=english: las referencias siguen la convención IEEE en inglés
% ("vol.", "no.", "pp.", "Accessed:", "[Online]. Available:"), como en la versión GDocs.
\usepackage[backend=biber,bibstyle=ieee,citestyle=numeric-comp,sorting=none,dashed=false,urldate=long,language=english,autolang=other]{biblatex}
\addbibresource{referencias.bib}
\DeclareSourcemap{\maps[datatype=bibtex]{\map[overwrite=false]{\step[fieldset=langid, fieldvalue=english]}}}
% english.lbx no activa \uspunctuation (coma dentro de las comillas de
% titulo, convencion IEEE); mapear a american.lbx, que si lo hace.
\DeclareLanguageMapping{english}{american}
% @online sin fecha: evitar el "()" vacío
\usepackage{xpatch}
\xpatchbibdriver{online}{\printtext[parens]{\usebibmacro{date}}}{\iffieldundef{year}{}{\printtext[parens]{\usebibmacro{date}}}}{}{}
\DeclareFieldFormat{sentencecase}{#1}   % respetar mayúsculas de los títulos tal como están en el .bib
\renewcommand*{\bibfont}{\small}
\appto{\bibsetup}{\raggedright}    % evita URLs "estiradas" en líneas justificadas
\setcounter{biburlnumpenalty}{9000}   % permite cortar URLs largas
\setcounter{biburllcpenalty}{7000}
\setcounter{biburlucpenalty}{8000}

% ---------- Hipervínculos (siempre al final) ----------
\usepackage[hidelinks,unicode]{hyperref}

% =====================================================================
%  Macros del proyecto
% =====================================================================
% Marcadores de trabajo pendiente — se ven en rojo/naranja en el PDF y se
% buscan fácil con:  grep -rn "pendiente\|confirmar\|revisar" capitulos/
% Para la versión final: \renewcommand{\pendiente}[1]{} etc. (o resolverlos).
\newcommand{\pendiente}[1]{\textcolor{red}{\textbf{[PENDIENTE:} #1\textbf{]}}}
\newcommand{\confirmar}[1]{\textcolor{orange!85!black}{\textbf{[CONFIRMAR:} #1\textbf{]}}}
\newcommand{\revisar}[1]{\textcolor{Purple}{\textbf{[REVISAR:} #1\textbf{]}}}
% Cambios de la última ronda de correcciones: en el PDF no se ven; en el Word
% (herramientas/word/resaltar.py) salen resaltados en amarillo, como pidió la
% Dra. Kruse (23 oct. 2026). Para la versión final se pueden borrar las marcas.
\newcommand{\resaltar}[1]{#1}

% Figura con respaldo: si el archivo no está, dibuja un recuadro con el nombre.
% Uso: \incluirfigura[width=0.8\linewidth]{Fig4_esquema_SLR.png}
\newcommand{\incluirfigura}[2][width=\linewidth]{%
  \IfFileExists{figuras/#2}%
    {\includegraphics[#1]{#2}}%
    {\fbox{\parbox[c][4cm][c]{0.8\linewidth}{\centering\small\ttfamily
      FALTA ARCHIVO: figuras/\detokenize{#2}}}}}

% Espacio extra en el índice para "Anexo A." (se llama desde 9_anexos.tex)
\makeatletter
\newcommand{\anexotocspacing}{\renewcommand*\l@section{\@dottedtocline{1}{1.5em}{5.2em}}}
% Índices de figuras y de tablas: más ancho para los números de dos cifras
% ("5.10.", "5.17."), que quedaban pegados al título (21 oct. 2026)
\renewcommand*\l@figure{\@dottedtocline{1}{1.5em}{2.9em}}
\renewcommand*\l@table{\@dottedtocline{1}{1.5em}{2.9em}}
% Índice general: menos espacio antes de cada capítulo (1,0em -> 0,55em) para que
% la última línea ("Referencias bibliográficas") no quede sola en una página
\patchcmd{\l@chapter}{\vskip 1.0em}{\vskip 0.55em}{}{\PackageWarning{tesis}{no se pudo ajustar el indice general}}
\makeatother

% Abreviaturas de uso frecuente
\newcommand{\stevia}{\textit{Stevia rebaudiana} Bertoni}
\newcommand{\Srebaudiana}{\textit{S.~rebaudiana}}
\newcommand{\anymaze}{AnyMaze\textsuperscript{\textcopyright}}

\usepackage{xr}
\externaldocument{build/main}

\begin{document}
\pagenumbering{arabic}

\begin{center}
  {\Large\textbf{Efectos neurobiológicos a largo plazo de la exposición a \textit{Stevia rebaudiana} durante la etapa juvenil y la adultez}}\\[0.8em]
  {\large Introducción, Estado del arte, Marco teórico, Perspectivas futuras y Conclusiones}\\[0.6em]
  Thiago Sayegh --- IByME-CONICET\\[0.3em]
  \textit{Versión para revisión de la Dra. M.\,S. Kruse --- 22 de octubre de 2026}
\end{center}

\vspace{0.8em}
\noindent\textbf{Cómo leer este documento.} Son las secciones 1, 2, 3, 7 y 8 de la tesis, seguidas de su propia lista de referencias (numeradas por orden de aparición, por lo que los números difieren de los de la tesis completa). Las referencias a otras secciones conservan la numeración de la tesis.

\clearpage
\setcounter{chapter}{0}
% =====================================================================
\chapter{Introducción}\label{cap:introduccion}
% Plantilla ITBA: máximo 3 carillas.
% 24 oct. 2026: texto reescrito por la Dra. Kruse (docx "tesina 20261024", sin
% control de cambios). Los párrafos iniciales siguen la introducción de Tsan et al.
% 2022 con sus referencias (comentario C7). \resaltarb{} = correcciones de Thiago
% sobre el texto de Sol en esta ronda (en el Word, en celeste).
% =====================================================================

Los niños constituyen el grupo etario con mayor consumo de azúcar, ya que aproximadamente el 16\,\% de su ingesta calórica total proviene de azúcares añadidos, de los cuales casi el 40\,\% corresponde a bebidas azucaradas\resaltarb{ \cite{ervin2012sugar}}. Estudios realizados tanto en humanos como en modelos de roedores han demostrado que el consumo excesivo de azúcar durante las etapas tempranas de la vida afecta negativamente el metabolismo de la glucosa\resaltarb{ \cite{vos2017sugars,cantoral2016ssb}}, el procesamiento sensorial y de la recompensa\resaltarb{ \cite{decock2016reward,jastreboff2016brain}}, la microbiota intestinal\resaltarb{ \cite{noble2021microbial,noble2017microbiome}} y las capacidades neurocognitivas\resaltarb{ \cite{cohen2018cognition,hsu2015sucrose,noble2019early,tsan2021western}}.

Con el objetivo de contrarrestar estos efectos adversos, los niños, al igual que los adultos, consumen cada vez más alimentos y bebidas que contienen edulcorantes bajos en calorías como sustitutos del azúcar\resaltarb{ \cite{sylvetsky2017lcs,sylvetsky2019pregnant}}. De hecho, el consumo de estos edulcorantes en la población infantil aumentó aproximadamente un 200\,\% entre 1999 y 2012\resaltarb{ \cite{sylvetsky2017lcs}}.

Aunque la sustitución de alimentos y bebidas con alto contenido de azúcar por productos que contienen edulcorantes bajos en calorías se considera una estrategia potencialmente beneficiosa para el control del peso corporal, la evidencia disponible sobre sus efectos en niños y adolescentes es contradictoria\resaltarb{ \cite{blum2005beverage,forshee2003beverage,striegel2006beverage,berkey2004beverages,sylvetsky2019lcsb,ludwig2001ssb,laska2012dietary}}.

En este contexto, los estudios mecanísticos realizados en modelos de roedores contribuyen a comprender los efectos a largo plazo del consumo habitual de edulcorantes durante distintas etapas de la vida. El presente trabajo busca evaluar si el consumo de stevia durante períodos críticos del desarrollo produce alteraciones neurobiológicas y conductuales en ratas macho Sprague-Dawley, con el fin de aportar evidencia sobre su seguridad. Asimismo, incorpora, como aporte desde la bioingeniería, un análisis multivariado de las métricas conductuales mediante técnicas de aprendizaje automático.

% ---------------------------------------------------------------------
\section{Planteamiento del problema}\label{sec:planteamiento}

El consumo de azúcar en Argentina constituye un problema de salud pública de primera magnitud: el argentino promedio ingiere el equivalente a 35 cucharaditas de azúcar por día \cite{fad2026bebidas}, los niños y adolescentes consumen en promedio 127 litros de bebidas azucaradas por año \cite{ennys2019}, y estas bebidas son responsables de más de 4.400 muertes anuales en el país, ubicándolo en el tercer lugar mundial en nuevos casos de diabetes tipo~2 atribuibles a su consumo \cite{laracastor2025ssb}. Frente a este panorama, el mercado de stevia crece sostenidamente (USD~840 millones en 2024) \cite{mordor2024stevia}, y Argentina es uno de sus principales mercados en América Latina \cite{ie2025latam,ie2025argentina}.

Sin embargo, la percepción pública de que la stevia es una alternativa inocua al azúcar no está suficientemente respaldada por evidencia científica en etapas críticas del desarrollo. La literatura ha documentado que el consumo temprano de sacarosa induce efectos neurobiológicos persistentes en la adultez: déficits en la memoria dependiente del hipocampo, alteraciones en la neurogénesis del giro dentado\resaltarb{ (GD)} y cambios conductuales relacionados con la ansiedad \cite{reichelt2016sucrose,coirini2022sucrose}. Trabajos del IByME-CONICET demostraron que la exposición temprana a stevia en ratas hembras induce alteraciones persistentes en la fertilidad, la descendencia y el comportamiento \cite{kruse2025stevia}. No obstante, los efectos a largo plazo sobre la conducta tipo ansiosa, la memoria de reconocimiento, la separación de patrones espaciales y los marcadores de neurogénesis en ratas macho aún no han sido caracterizados.

La pregunta central que guía este trabajo es si el consumo de stevia durante etapas tempranas del desarrollo induce alteraciones conductuales y neurobiológicas que persisten hasta la edad adulta, si estas son comparables con las producidas por el consumo de sacarosa y en qué medida dependen de la etapa del desarrollo en la que ocurre la exposición.

Para abordar estos interrogantes, se utilizaron ratas macho Sprague-Dawley expuestas a los endulzantes durante dos períodos del desarrollo distintos: una etapa juvenil, comprendida entre los días \resaltarb{postnatales} (PD) 25 y 50, que correspondería a la transición infanto-\resaltarb{adolescente} en la rata\resaltarb{ \cite{spear2000adolescent}}, y una etapa adulta, comprendida entre PD75 y \resaltarb{PD100}. Se evaluaron las conductas de tipo ansioso y el desempeño en pruebas de memoria, así como la expresión del antígeno nuclear de proliferación celular (PCNA) y de doblecortina (DCX), marcadores asociados con distintas etapas del proceso de neurogénesis, en el \resaltarb{GD} del hipocampo. Este abordaje permitió comparar los efectos del consumo de stevia y sacarosa en función de la edad de exposición y explorar su posible relación con alteraciones persistentes en la neurogénesis hipocampal.

% ---------------------------------------------------------------------
\section{Justificación}\label{sec:justificacion}

El elevado consumo de azúcar durante la infancia y la adolescencia constituye una preocupación para la salud pública, debido a sus efectos adversos sobre la salud metabólica y sus consecuencias sobre el desarrollo cerebral y las funciones cognitivas. En este contexto, los edulcorantes bajos en calorías, como la stevia, han adquirido creciente relevancia como alternativas al azúcar. Sin embargo, aún es necesario profundizar en sus posibles efectos sobre el sistema nervioso central, especialmente durante etapas críticas del desarrollo.

La etapa juvenil constituye un período de especial interés debido a los procesos de maduración y reorganización de los circuitos neuronales que tienen lugar durante esta ventana del desarrollo. En consecuencia, la exposición a determinados componentes de la dieta podría producir modificaciones neurobiológicas y conductuales que persistan hasta la edad adulta.

En este marco, resulta relevante investigar si el consumo de stevia durante la etapa juvenil produce alteraciones en la memoria, las conductas de tipo ansioso y los procesos de neurogénesis hipocampal, y determinar si estos efectos difieren de los observados cuando la exposición ocurre en la edad adulta. Asimismo, la comparación con la sacarosa, cuyos efectos adversos sobre determinados procesos cognitivos y neurobiológicos han sido documentados previamente, permite contextualizar la magnitud y las características de las posibles alteraciones asociadas al consumo de stevia.

De esta manera, el presente estudio busca aportar evidencia experimental sobre los efectos neurobiológicos y conductuales del consumo de stevia en distintas etapas del desarrollo, contribuyendo a una evaluación más integral de su seguridad como alternativa al azúcar.

La pertinencia en el marco de la carrera de Bioingeniería es doble. Por un lado, el trabajo implica técnicas experimentales ---pruebas conductuales estandarizadas, inmunocitoquímica seguida de análisis por \resaltarb{microscopía} confocal--- que requieren precisión técnica y criterio científico. Por otro, los objetivos de máxima incorporan herramientas de ingeniería de datos y aprendizaje automático ---PCA, clustering jerárquico y Random Forest--- cuya aplicación sobre datos conductuales constituye un aporte metodológico propio de la formación en ingeniería.

% ---------------------------------------------------------------------
\section{Objetivos}\label{sec:objetivos}

\subsection{Objetivo general}

Evaluar los efectos neuroconductuales a largo plazo del consumo de \stevia{} en ratas macho Sprague-Dawley, en comparación con los inducidos por el consumo de sacarosa, y determinar en qué medida estos efectos dependen de la etapa del desarrollo en la que ocurre la exposición (juvenil o adulta).

\subsection{Objetivos de mínima}\label{sec:obj-minima}

\begin{enumerate}[label=\arabic*.]
  \item Registrar el consumo de stevia/sacarosa y agua (día por medio) en los dos grupos etarios (\resaltarb{PD25--50 y PD75--100}) durante el tratamiento.
  \item Evaluar la conducta tipo ansiosa mediante la prueba de \textit{open field} (OF) al finalizar el tratamiento, contrastando los grupos juvenil y adulto y las condiciones stevia, sacarosa y control (6 grupos experimentales, en total).
  \item Cuantificar la memoria de reconocimiento de objetos a corto y largo plazo y la memoria consolidada (NOR: T2, T3 y T4) mediante la tasa de exploración.
  \item Evaluar la separación de patrones espaciales mediante la prueba SLR (configuraciones d-SLR y s-SLR).
  \item Evaluar si la neurogénesis se encuentra alterada por el tratamiento. Para ello se medirá la expresión de los progenitores neuronales (PCNA) y de neuronas inmaduras (DCX) mediante inmunocitoquímica \resaltarb{seguida de microscopía} confocal.
\end{enumerate}

\subsection{Objetivos de máxima}\label{sec:obj-maxima}

\begin{enumerate}[label=\arabic*.]
  \item Aplicar análisis de componentes principales (PCA) y clustering jerárquico sobre el conjunto de métricas conductuales extraídas con \anymaze{} para identificar perfiles conductuales diferenciados entre grupos sin asumir separación a priori.
  \item Desarrollar un modelo de clasificación supervisada (Random Forest) entrenado con las variables conductuales de \anymaze{} como predictores, con el objetivo de determinar la capacidad de dichas métricas para predecir el grupo de exposición de cada animal, entre los seis grupos experimentales resultantes del cruce de la ventana etaria (juvenil, adulto) y la condición dietaria (control, sacarosa, stevia), y cuantificar la importancia relativa de cada variable conductual.
\end{enumerate}

% ---------------------------------------------------------------------
\section{Hipótesis}\label{sec:hipotesis}

Se postula que el consumo de \stevia{} durante etapas tempranas del desarrollo no reproducirá los efectos adversos sobre el comportamiento y la neurogénesis hipocampal documentados para la sacarosa \cite{reichelt2016sucrose,coirini2022sucrose}. Por el contrario, en virtud de las propiedades hipoglucemiantes, antiinflamatorias y neuroprotectoras de los glucósidos de esteviol, se espera que la stevia actúe como un compuesto inocuo o eventualmente beneficioso sobre los parámetros conductuales y los marcadores de neurogénesis evaluados \cite{kruse2025stevia,pozdnyakova2025stevia,prata2017glycosides}. En consecuencia, los animales expuestos a stevia ---tanto juveniles como adultos--- mostrarán un perfil conductual y neurobiológico comparable al de los controles, o incluso más favorable, en contraste con los déficits esperables en los grupos expuestos a sacarosa. El efecto diferencial según la ventana etaria permitirá determinar si las alteraciones halladas dependen de la edad de exposición, como ocurre con el consumo de sacarosa.

% =====================================================================
\chapter{Estado del arte}\label{cap:estado-arte}
% 24 oct. 2026 (correcciones de la Dra. Kruse): sin el párrafo inicial; la
% descripción de la literatura va en presente; nuevo título de §2.3.
% =====================================================================

% ---------------------------------------------------------------------
\section{Edulcorantes no calóricos y \textit{Stevia rebaudiana}}\label{sec:stevia-evidencia}

\resaltarb{Los edulcorantes no calóricos no son metabólicamente inertes: modifican la respuesta de las regiones cerebrales vinculadas a la recompensa y al procesamiento del sabor dulce} \cite{yunker2020nns}\resaltarb{. En ratas, el consumo diario de edulcorantes bajos en calorías (acesulfamo de potasio, sacarina o stevia) desde la etapa juvenil altera la conducta motivada por el azúcar y la memoria dependiente del hipocampo; sin embargo, la memoria se evaluó mientras los animales todavía consumían los edulcorantes, y el deterioro se observó al analizar los tres edulcorantes en conjunto, no la stevia por separado} \cite{tsan2022lcs}. La stevia, en particular, ocupa una posición ambigua. Por un lado, se le atribuyen propiedades neuroprotectoras: reduce la peroxidación lipídica cerebral en ratas alimentadas con una dieta alta en grasas y sacarosa \cite{pozdnyakova2025stevia}, sus glucósidos presentan actividad antioxidante e insulino-mimética \cite{prata2017glycosides}, el esteviósido \resaltarb{tiene} un efecto antiamnésico en ratas tratadas con escopolamina \cite{singh2010stevioside}, y la stevia \resaltarb{modula} la plasticidad sináptica en un modelo de desorden metabólico \cite{chavushyan2017stevia} y \resaltarb{revierte} la ansiedad y el deterioro de la memoria en ratas diabéticas \cite{khakpai2023stevia}. Por otro lado, su consumo \resaltarb{activa} el sistema de recompensa \cite{salinasvelarde2021fosb}, y \resaltarb{estudios previos del laboratorio muestran que} la exposición temprana a un extracto acuoso de stevia \resaltarb{altera} la fertilidad y la descendencia de ratas hembra, a la vez que \resaltarb{reduce} su ansiedad \cite{kruse2025stevia}. En ratas macho, el consumo temprano del mismo extracto no \resaltarb{modifica} la conducta tipo ansiosa ni la memoria cuando se \resaltarb{evalúa mientras dura el consumo} \cite{coirini2023stevia,marchena2025stevia}; no se dispone, sin embargo, de estudios que evalúen sus efectos a largo plazo, tras el retiro del edulcorante, con una batería conductual e histológica en las ventanas etarias juvenil y adulta, vacío que este proyecto aborda.

% ---------------------------------------------------------------------
\section{Consumo temprano de azúcar, memoria y neurogénesis hipocampal}\label{sec:azucar-neurogenesis}

La sacarosa aporta el marco de referencia contra el cual contrastar la stevia. \resaltarb{Estudios previos} del laboratorio muestran que la exposición juvenil a sacarosa al 10\,\% \resaltarb{aumenta} la conducta tipo ansiosa \cite{kruse2019sucrose} y altera la memoria de reconocimiento de objetos en la adultez, sin efectos cuando la exposición \resaltarb{ocurre} en animales adultos \cite{coirini2022sucrose}. Otros trabajos \resaltarb{describen} deterioros de la memoria dependiente del hipocampo tras el consumo de azúcar \cite{hsu2015sucrose,xu2018sucrose}, que \resaltarb{persisten} en la adultez cuando la exposición se \resaltarb{restringe} a la adolescencia \cite{noble2019early,reichelt2015adolescent}, y el acceso diario a sacarosa \resaltarb{altera} la discriminación espacial cuando los objetos \resaltarb{están} próximos, una tarea que depende de la separación de patrones, a la vez que \resaltarb{reduce} la proliferación en el \resaltarb{GD} \cite{reichelt2016sucrose}. Entre los mecanismos propuestos se cuentan la alteración de la microbiota intestinal \cite{noble2021microbial} y la menor expresión de genes hipocampales vinculados a la neurogénesis \cite{maniam2016sugar}. Los resultados dependen, no obstante, de la ventana de exposición: el consumo de sacarosa iniciado en la adolescencia tardía no \resaltarb{induce} conducta tipo ansiosa y \resaltarb{aumenta} la proliferación en la zona subgranular \cite{sanchezhuerta2022sucrose}, lo que respalda la elección de dos ventanas etarias contrastantes (PD25 y PD75).

% ---------------------------------------------------------------------
\section{\resaltarb{Aprendizaje automático aplicado al análisis del comportamiento}}\label{sec:ml-fenotipado}

El análisis conductual tradicional se apoya en medidas univariadas que, examinadas de forma aislada, omiten la estructura conjunta del comportamiento \cite{geros2020tracking}. El aprendizaje automático se ha aplicado a las métricas que los sistemas de seguimiento registran en cada animal: modelos de Random Forest discriminan fenotipos asociados a disfunción dopaminérgica o serotoninérgica e identifican las variables de mayor peso \cite{lukovikov2026cylinder}, y otros clasificadores \resaltarb{diferencian} cepas de ratas a partir de sus perfiles conductuales \cite{geros2020tracking}. En ambos casos, los grupos \resaltarb{difieren} por manipulaciones farmacológicas o genéticas de gran efecto; su aplicación a intervenciones nutricionales, de efectos más sutiles, permanece escasamente explorada.

% =====================================================================
\chapter{Marco teórico}\label{cap:marco-teorico}
% Plantilla ITBA: sustento formal que valida las decisiones de diseño;
% evitar el "efecto enciclopedia".
% 24 oct. 2026: §3.1, §3.2 y §3.3 con los textos de la Dra. Kruse. Sus números de
% cita venían de otros trabajos: [42–45] de Coirini et al. 2022 (Biomedicines),
% [33–35] del borrador de Marchena et al. y [56–62] de un borrador del laboratorio
% sin lista de referencias (se citan las fuentes primarias, verificadas en PubMed).
% =====================================================================

% ---------------------------------------------------------------------
\section{Ventanas críticas del neurodesarrollo}\label{sec:ventanas}

El estudio de los efectos del consumo de bebidas azucaradas durante las etapas tempranas de la vida reviste especial interés, considerando el elevado consumo de estos productos entre niños y adolescentes y su incremento registrado en las últimas décadas \cite{guthrie2000sweeteners,frary2004children,guelinckx2015water,wang2008ssb}. Esta etapa del desarrollo se caracteriza por una intensa maduración y reorganización del sistema nervioso central, acompañada de una elevada plasticidad neuronal y de la adquisición de nuevos conocimientos y habilidades \cite{sisk2004puberty,spear2000adolescent}. Durante este período tienen lugar procesos fundamentales, como la maduración de los circuitos neuronales, el refinamiento de las conexiones sinápticas y el desarrollo progresivo de las funciones cognitivas y emocionales. En consecuencia, la exposición a determinados factores ambientales y nutricionales durante estas ventanas de especial vulnerabilidad podría interferir con dichos procesos y producir alteraciones neurobiológicas y conductuales que persistan hasta la edad adulta.

En este contexto, los modelos experimentales de consumo de bebidas azucaradas durante la etapa juvenil permiten investigar si la exposición temprana a estos productos afecta el desarrollo y funcionamiento del sistema nervioso central, particularmente en relación con los procesos de aprendizaje, memoria y regulación emocional. Asimismo, permiten evaluar si las consecuencias de dicha exposición dependen de la etapa del desarrollo en la que ocurre el consumo y si sus efectos se mantienen una vez finalizada la exposición.

Por otra parte, dado que los edulcorantes bajos en calorías, como la stevia, se utilizan como alternativas al azúcar, resulta igualmente relevante investigar sus posibles efectos sobre el sistema nervioso central durante estas etapas del desarrollo. Si bien su utilización permite reducir el aporte de azúcares de la dieta, aún es necesario profundizar en las posibles consecuencias neurobiológicas y conductuales de su consumo prolongado. Por ello, comparar los efectos de la exposición a stevia y sacarosa durante las etapas juvenil y adulta constituye una estrategia experimental de interés para determinar si ambos tratamientos producen alteraciones persistentes y establecer en qué medida la edad de exposición condiciona la respuesta del sistema nervioso. De esta manera, el presente estudio busca aportar evidencia sobre los posibles efectos a largo plazo del consumo de stevia y contribuir a una evaluación más integral de su seguridad como alternativa al azúcar, especialmente durante períodos críticos del neurodesarrollo.

% ---------------------------------------------------------------------
\section{Fundamento del esquema de exposición}\label{sec:bases}

La inclusión de tres grupos experimentales responde a una estrategia de comparación que permite contextualizar los efectos de los distintos tratamientos. En particular, la administración de sacarosa al 10\,\% p/v se utiliza como modelo experimental del consumo de bebidas azucaradas, cuyos efectos adversos sobre la memoria, las conductas de tipo ansioso y la neurogénesis han sido documentados previamente \cite{reichelt2016sucrose,coirini2022sucrose,kruse2019sucrose}. La sacarosa es hiperglucemiante, y su consumo elevado se asocia con procesos inflamatorios y alteraciones metabólicas \cite{anton2010stevia,vos2017sugars,hsu2015sucrose}. El grupo control corresponde a los animales que consumen solamente agua durante todo el período experimental. La stevia, administrada como extracto acuoso de hojas al 4\,\% v/v (\S\ref{sec:soluciones}), aporta sabor dulce prácticamente sin aportar calorías \cite{chatsudthipong2009stevioside}. A diferencia de la sacarosa, la stevia presenta efectos hipoglucemiantes y antioxidantes, por lo que es considerada una alternativa alimentaria para personas con trastornos metabólicos \cite{anton2010stevia,chatsudthipong2009stevioside,gregersen2004stevioside,prata2017glycosides}. Para este estudio, se optó por utilizar un extracto acuoso natural de \textit{Stevia rebaudiana}, con el propósito de evitar la presencia de aditivos u otros edulcorantes que suelen encontrarse en las formulaciones comerciales y que podrían interferir con los resultados experimentales. El extracto se preparó siguiendo un método previamente descrito \cite{kruse2025stevia,sharma2009stevia}, seleccionado por su elevado rendimiento en glucósidos de esteviol \cite{celaya2022infusions} y por su similitud con los procedimientos de preparación difundidos en portales de salud dirigidos a la comunidad. La composición química del extracto utilizado se detalla en Materiales y métodos (Tabla~\ref{tab:extracto}).

% ---------------------------------------------------------------------
\section{Neurogénesis hipocampal y separación de patrones}\label{sec:cascada}

En el GD, la formación de nuevas neuronas ocurre en etapas sucesivas: las células madre y los progenitores en proliferación (PCNA$^{+}$) dan lugar a neuroblastos y neuronas inmaduras (DCX$^{+}$), que maduran en neuronas granulares que se integran al circuito y expresan el antígeno nuclear neuronal (NeuN$^{+}$; Figura~\ref{fig:cascada-neurogenica}) \cite{kempermann2015neurogenesis}.

\begin{figure}[htbp]
  \centering
  \incluirfigura[width=0.7\linewidth]{Fig3_cascada_neurogenica.jpg}
  \caption[Esquema de la cascada neurogénica del GD]{Esquema de la cascada neurogénica del GD y de los marcadores que indexan cada etapa: PCNA$^{+}$ (células madre y progenitores en proliferación), DCX$^{+}$ (neuroblastos y neuronas inmaduras) y NeuN$^{+}$ (neuronas granulares maduras). Elaboración propia a partir de Kempermann et al.~\cite{kempermann2015neurogenesis}.}
  \label{fig:cascada-neurogenica}
\end{figure}

PCNA es un marcador de proliferación de células progenitoras neurales \cite{eisch2007beyond,mandyam2007precursor}, que se expresa durante todas las fases activas del ciclo celular y durante un breve período después de que las células adquieren un estado postmitótico \cite{reichelt2016sucrose}. Por su parte, DCX es una proteína asociada a los microtúbulos que se expresa en células precursoras neuronales en migración \cite{gleeson1999dcx,francis1999dcx} y se utiliza como marcador de neuronas inmaduras generadas durante la neurogénesis adulta en el GD \cite{brown2003dcx,couillard2005dcx}.

Se ha propuesto que la neurogénesis adulta en el GD constituye un sustrato neuronal para determinados aspectos de la formación de la memoria, particularmente para la separación de patrones (\textit{pattern separation}), un proceso que permite distinguir y codificar patrones similares de actividad neuronal como representaciones diferentes durante la codificación y el almacenamiento de la información \cite{reichelt2016sucrose,reichelt2021slr}. En este sentido, la neurogénesis adulta proporcionaría una funcionalidad adicional al hipocampo, especialmente cuando aumenta la demanda cognitiva \cite{reichelt2016sucrose,reichelt2021slr}.

% ---------------------------------------------------------------------
\section{Fundamento del abordaje multivariado}\label{sec:multivariado}

Las métricas que \anymaze{} registra en cada animal ---en la SLR, el tiempo de exploración de cada objeto en cada sesión, la exploración total y el índice de discriminación (DI)--- están correlacionadas entre sí. Analizarlas una por una infla la tasa de error por familia de contrastes (FWER),
\begin{equation}
  \mathrm{FWER} = 1 - (1-\alpha)^{m},
  \label{eq:fwer}
\end{equation}
donde $m$ es el número de contrastes, y descarta su estructura conjunta: una diferencia entre grupos puede residir en una combinación de variables sin que ninguna alcance significancia por separado.

El PCA resume las variables estandarizadas en combinaciones no correlacionadas mediante la descomposición espectral de su matriz de correlación \cite{jolliffe2016pca},
\begin{equation}
  \mathbf{R} = \mathbf{V}\boldsymbol{\Lambda}\mathbf{V}^{\mathsf{T}},
  \label{eq:pca}
\end{equation}
cuyos autovectores definen las componentes y cuyos autovalores $\lambda_k$ cuantifican la varianza que explica cada una. Sobre las mismas variables estandarizadas, el clustering jerárquico con el criterio de Ward, que minimiza el aumento de la varianza dentro de los grupos al fusionarlos \cite{ward1963hierarchical}, evalúa si los animales se agrupan espontáneamente según su exposición (objetivo de máxima 1); el coeficiente de silueta indica el número de grupos naturales, y el índice de Rand ajustado (ARI), su coincidencia con los grupos experimentales (0: azar; 1: coincidencia perfecta).

El objetivo de máxima 2 ---determinar si las métricas predicen el grupo de cada animal y con qué peso--- es un problema de clasificación supervisada, para el cual el Random Forest, un conjunto de árboles de decisión entrenados sobre muestras \textit{bootstrap}, resulta adecuado por tolerar relaciones no lineales y variables correlacionadas \cite{breiman2001rf}. Con los tamaños muestrales de la experimentación animal \cite{charan2013sample}, su desempeño se estima mediante validación cruzada (el modelo se evalúa con animales que no usó para entrenarse) y se compara con el obtenido al mezclar las etiquetas al azar (prueba de permutación); la importancia de cada variable es la caída del desempeño al permutar sus valores. Como el clasificador aprovecha cualquier diferencia sistemática entre grupos, deben excluirse las variables que identifiquen la cohorte de origen y no el tratamiento.

\setcounter{chapter}{6}
% =====================================================================
\chapter{Perspectivas futuras}\label{cap:perspectivas}
% Plantilla ITBA: opcional, máximo 1 carilla. Puede incluir objetivos de
% máxima no alcanzados. Borrador del 21 oct. 2026: incluye, como trabajo a
% futuro, las limitaciones del diseño que la Dra. Kruse sacó de la Discusión
% (a confirmar con ella).
% =====================================================================

Los resultados de este trabajo abren distintas líneas para continuarlo:

\begin{itemize}
  \item \textbf{Neurogénesis en los grupos expuestos a stevia.} Cuantificar PCNA y DCX en los animales expuestos a stevia permitirá determinar si la preservación de la separación de patrones espaciales se acompaña de una neurogénesis conservada. Conviene, además, registrar el animal de origen de cada corte, de modo que el análisis pueda realizarse por animal y no por corte.
  \item \textbf{Mecanismos.} Evaluar si las modificaciones en la expresión de calretinina y parvalbúmina descriptas tras el consumo agudo de stevia persisten a largo plazo y se relacionan con las alteraciones de la memoria de reconocimiento, y estudiar el papel de la abstinencia de la bebida endulzada en los cambios de la conducta tipo ansiosa.
  \item \textbf{Hembras.} Extender el estudio a ratas hembra, en las que el mismo extracto redujo la conducta tipo ansiosa y alteró la fertilidad, para establecer si esos efectos persisten a largo plazo y si dependen de la edad de exposición.
  \item \textbf{Dosis.} Dado que la ingesta superó los valores diarios de referencia, evaluar concentraciones menores del extracto, cercanas a la ingesta diaria admisible de glucósidos de esteviol.
  \item \textbf{Diseño de las cohortes.} Incluir un grupo control en cada cohorte experimental y ampliar el número de animales expuestos a stevia en el \textit{open field}.
  \item \textbf{Análisis multivariado.} Registrar las métricas de \anymaze{} identificadas por animal también en el \textit{open field} y en el NOR, para extender el PCA, el clustering y el Random Forest a toda la batería conductual e integrar en un mismo análisis las distintas pruebas, junto con el peso corporal.
\end{itemize}

% =====================================================================
\chapter{Conclusiones}\label{cap:conclusiones}
% Plantilla ITBA: máximo 2 carillas; indicar si se cumplieron los
% objetivos y si la hipótesis fue comprobada o refutada.
% Borrador del 21 oct. 2026 (pendiente de revisión de la Dra. Kruse): una
% conclusión por asunto, estado de cada objetivo y de la hipótesis. No agrega
% resultados nuevos: resume los de la sección 5 con el enfoque de la Discusión.
% =====================================================================

En este trabajo se evaluaron los efectos a largo plazo del consumo de una bebida endulzada con un extracto de stevia durante la etapa juvenil o la adultez, y se los comparó con los de la sacarosa, utilizada como control positivo por sus efectos ya descriptos. De los resultados obtenidos se desprenden las siguientes conclusiones.

\textbf{Consumo.} Los animales juveniles consumieron, en relación con su peso corporal, más del doble de solución endulzada que los adultos, sin una mayor preferencia por el sabor dulce, y la ingesta de ambos \resaltar{endulzantes} superó los valores diarios de referencia.

\textbf{Conducta tipo ansiosa.} El consumo de stevia, en cualquiera de las dos edades, se asoció a largo plazo con una mayor exploración de las zonas centrales del \textit{open field}, sin cambios en la actividad locomotora, es decir, con un perfil de tipo ansiolítico. La sacarosa, en cambio, aumentó la conducta tipo ansiosa solo cuando se consumió durante la etapa juvenil.

\textbf{Memoria de reconocimiento.} La stevia alteró la memoria de reconocimiento de objetos de manera dependiente de la edad: a largo plazo en los animales expuestos durante la etapa juvenil y a corto plazo en los expuestos durante la adultez, sin afectar la memoria consolidada. La sacarosa consumida durante la etapa juvenil alteró los tres tipos de memoria evaluados.

\textbf{Separación de patrones espaciales y neurogénesis.} La stevia no alteró la separación de patrones espaciales en ninguna de las dos edades. La sacarosa la deterioró en ambas edades en la configuración de mayor demanda, y también en la de menor demanda cuando se consumió durante la etapa juvenil, y redujo en ambas edades la proliferación celular y el número de neuroblastos en proliferación en el \resaltarb{GD}: dos expresiones de un mismo efecto sobre la neurogénesis hipocampal.

\textbf{Análisis multivariado.} Las métricas de la prueba SLR se organizaron en dos dimensiones casi independientes, la exploración de los objetos y la discriminación de la posición novedosa. Los animales expuestos a stevia se ubicaron junto a los controles en la dimensión de la discriminación, y el clasificador distinguió a cada \resaltar{endulzante} del control por encima del azar, aunque ningún grupo constituyó un perfil conductual separado del resto.

\section*{Cumplimiento de los objetivos}

Se cumplieron los objetivos de mínima referidos al registro del consumo, a la conducta tipo ansiosa, a la memoria de reconocimiento y a la separación de patrones espaciales. \resaltar{La} cuantificación de PCNA y DCX se completó en los grupos control y sacarosa, pero no en los grupos expuestos a stevia. Los objetivos de máxima se cumplieron sobre la prueba SLR: el PCA y el clustering de Ward mostraron que los animales no se agrupan espontáneamente según su exposición, y el Random Forest predijo el grupo de cada animal por encima del azar, con el índice de discriminación como la variable de mayor importancia.

\section*{Hipótesis}

La hipótesis se cumplió solo en parte. La stevia no reprodujo los efectos adversos de la sacarosa sobre la memoria consolidada y la separación de patrones espaciales, pero no resultó inocua: produjo, en ambas edades, un perfil de tipo ansiolítico y alteraciones de la memoria de reconocimiento, de menor magnitud que las de la sacarosa. Su efecto sobre la neurogénesis hipocampal queda por determinar.

En conclusión, el consumo de stevia en etapas tempranas de la vida o en la adultez no está exento de efectos a largo plazo sobre el comportamiento, aunque estos son más sutiles que los de la sacarosa. Estos resultados indican que la seguridad de la stevia como alternativa al azúcar, especialmente durante el desarrollo, requiere una evaluación más profunda.


\chapter*{Referencias bibliográficas}
\printbibliography[heading=none]

\end{document}
